\section{Mein Beitrag für das Schulganze}
\subsection{Was ist mein aktueller Beitrag für das Schulganze (Fachschaft, Arbeitsgruppen,
	Schilf, ...)?}
\begin{itemize}
	\item Coaching für elektronisches Prüfen für andere Fachschaften
	\item Gemeinsames Prüfen stärken (e.g. Jahrgangsprüfung Mathematik)
	\item Beitrag MINT-Tag und MINT-Tagung (Vortrag mit gutem Feedback)
	\item randomisierte Übungsaufgaben auf Moodle (Projekt)
	\item Erstellung von GitHub-Repository für gemeinsame Unterlagen sowie dessen Verwaltung (Administration). Unsere aktuellen Unterlagen finden sich hier: \url{https://gofile.me/6YFll/vyDysf7zw}
	\item Einrichtung und Unterhalt von Infrastruktur für die Fachschaft Informatik (später hoffentlich auch für die Fachschaft Mathematik).
	\item \LaTeX-Setup für die fortlaufende gemeinsame Erstellung und Verbesserung des Unterrichtsmaterials (gekoppelt mit GitHub)
	\item technische Unterstützung von Lehrpersonen (Beantwortung technischer Fragen)
	\item Anbieten von Praktikumsplätzen (Lehrdiplom Informatik)
\end{itemize}
\subsection{Wo habe ich Interesse, Ideen und Stärken, die ich für das Schulganze zur Verfügung stellen möchte?}
\begin{itemize}
	\item Aufbau und Dokumentation eines Workflows zur gemeinsamen Erstellung von Unterlagen, woran auch andere Fachschaften interessiert sein könnten
	\item Jahrgangsprüfung / gemeinsames Prüfen (siehe oben)
	\item Was würden andere Lehrpersonen gerne sehen? Wo besteht Unterstützungsbedarf?
\end{itemize}

\clearpage
\section{Unterricht}
\subsection{Was gelingt mir gut im Unterricht?}
\begin{itemize}
	\item SuS schätzen meine starke Identifikation mit dem Fach
	\item Einbettung des Inhalts in den grösseren Zusammenhang
	\item Es fällt mir recht leicht einen Bezug zu den SuS aufzubauen
	\item Klassenführung klar besser geworden
\end{itemize}
\subsection{Mit welchen Aspekten möchte ich mich vertieft auseinandersetzen?}
\begin{itemize}
	\item Feinschliff von Unterlagen
	\item Noch langer Weg um herauszufinden, welche Reihenfolge und Auswahl der Themen am besten funktioniert 
\end{itemize}
\subsection{Welche Weiterbildungsziele habe ich?}
\begin{itemize}
	\item Weiterbildung in Emmetten (Praktikumslehrperson): Montag, 9. Februar 2026 (Beginn um 9:30 Uhr) bis Mittwoch, 11. Februar 2026 (Abreise um 15:00 Uhr)
	\item Teilnahme am Swiss STACK User Meeting 2025 (Mai 2025) und halten eines Vortrags an diesem Meeting
	\item Fachliche Weiterbildungungen:\\
	\url{https://educ.ethz.ch/lernzentren/mint-lernzentrum/weiterbildungsangebote/fortbildungsangebote-im-fach-physik/Vom-Doppelspalt-zum-Quantencomputer.html} 
	\item Englisch C2-Zertifikat für Immersionsklassen
	\item Weiterführende Bearbeitung diverser Projekte
	\item Später eventuell Lehrdiplom Mathematik
\end{itemize}

\clearpage
\section{Meine Anliegen}
\subsection{Welche Anliegen habe ich bezüglich Schulentwicklung? Gibt es Bereiche, die
	stärker beachtet werden sollten?}
\begin{itemize}
	\item Möglichst geringe Schwankungen in der Anzahl der Lektionen zwischen HS und FS
\end{itemize}

\subsection{Wo wünsche ich Unterstützung durch die Schulleitung?}
\begin{itemize}
	\item Garantiertes Pensum (zur Zeit kann ich mein Pensum nur durch Abbau von Überzeit aufrechterhalten)
	\item Klare und vorausschauende Kommunikation: Wann finden Leistungsbeurteilungen statt / wie geht es weiter mit Stufenanstiegen
\end{itemize}

\subsection{Weitere Punkte?}
\begin{itemize}
	\item Beginner's Guide to Lee
	\item Viel schnellere Aktualisierung der LP-Porträts auf der Lee-Webseite
\end{itemize}